\chapter{Données point-in-time et biais}\label{ch:rs:point-in-time-data-and-the-biases}

Le 3 octobre 2008, le rapport américain sur l'emploi annonçait que l'économie avait perdu 159\,000 emplois en septembre. La base actuelle indique
451\,000. Entre les deux, le chiffre a été révisé six fois, dont trois lors des révisions annuelles de référence que
l'agence statistique publie chaque février. Un signal macroéconomique testé en backtest sur
la base actuelle se serait fondé sur un chiffre que personne n'avait en octobre 2008, et aurait vu une récession plus profonde et plus nette
que celle que quiconque voyait à l'époque. Ce chapitre traite de la différence entre ce qui est vrai aujourd'hui et ce qui était connu alors~:
le \omterm{def:rs:point-in-time-data-and-the-biases:lookahead}{biais d'anticipation} et les estampilles temporelles qui l'empêchent, le \omterm{def:rs:point-in-time-data-and-the-biases:survivorship}{biais de survie}, les révisions et les \omterm{def:rs:point-in-time-data-and-the-biases:vintage}{retraitements}, le stockage qui conserve
chaque millésime d'un fait, et la \omterm{def:rs:point-in-time-data-and-the-biases:asof}{jointure as-of} qui le lit. Ses données sont les millésimes en temps réel des créations d'emplois non agricoles américaines et un
marché d'actions simulé, avec radiations.

\section{Biais d'anticipation et sémantique des estampilles temporelles}

\begin{definition}[Biais d'anticipation, données point-in-time]\label{def:rs:point-in-time-data-and-the-biases:lookahead}
Un backtest souffre d'un \emph{biais d'anticipation}\index{biais d'anticipation} (look-ahead bias) lorsqu'une décision prise à l'instant $t$ utilise une information qui n'était pas
disponible à l'instant $t$. Les \emph{données point-in-time}\index{données point-in-time} sont des données stockées de façon que, pour tout instant passé, elles
restituent exactement ce qui était connu à cet instant~: les valeurs alors publiées, l'univers alors coté, les identifiants alors
utilisés.
\end{definition}

Le \omterm{def:rs:point-in-time-data-and-the-biases:lookahead}{biais d'anticipation} vient rarement d'un regard délibéré. Il provient d'une estampille temporelle qui signifie autre chose que ce que le code
suppose. Une barre quotidienne estampillée de sa date n'est connue qu'à la clôture, si bien qu'un signal calculé à partir d'elle ne peut pas
s'exécuter à la clôture de ce jour-là. Un chiffre de résultats estampillé de la fin de son trimestre comptable n'est connu que des semaines plus tard, à la date du dépôt (chapitre~11).
Un changement d'indice estampillé de sa date d'entrée en vigueur a été annoncé quelques jours plus tôt (le livre~1, chapitre~15), et
l'annonce est le moment où la transaction doit commencer. Une transaction estampillée par la bourse parvient à la société plus tard, et
l'heure d'arrivée est celle sur laquelle la société pouvait agir (livre~1, chapitre~28). Chaque fait exige donc deux instants.

\begin{definition}[Temps de validité, temps de connaissance, données bitemporelles]\label{def:rs:point-in-time-data-and-the-biases:bitemporal}
Le \emph{temps de validité}\index{temps de validité} (valid time) d'une valeur est l'instant qu'elle décrit~: septembre 2008 pour la variation de l'emploi de septembre, le
quatrième trimestre pour un chiffre trimestriel de résultats. Son \emph{temps de connaissance}\index{temps de connaissance} (knowledge time) est l'instant où elle est devenue accessible à
l'utilisateur~: la publication, le dépôt, la livraison par le fournisseur, l'arrivée du message. Des données sont \emph{bitemporelles}\index{bitemporelles} (bitemporal)
lorsque chaque valeur est stockée avec les deux, et qu'une nouvelle valeur pour un même temps de validité est ajoutée plutôt que d'écraser l'ancienne.
\end{definition}

\begin{omfigure}
\begin{tikzpicture}[font=\footnotesize,x=1.35cm,y=0.8cm]
\draw[->,>=Stealth] (0,0) -- (8.4,0) node[right] {temps de validité};
\draw[->,>=Stealth] (0,0) -- (0,5.4) node[above] {temps de connaissance};
\foreach \x/\lab in {1/août,2/sept.,3/oct.,4/nov.} \node[below] at (\x,0) {\lab};
\foreach \y/\lab in {1/oct.,2/nov.,3/déc.,4/{févr. ou après}} \node[left] at (0,\y) {\lab};
\fill[omDef] (1,1) circle (2.2pt) (1,2) circle (2.2pt) (1,3) circle (2.2pt) (1,4) circle (2.2pt);
\fill[omDef] (2,1) circle (2.2pt) (2,2) circle (2.2pt) (2,3) circle (2.2pt) (2,4) circle (2.2pt);
\fill[omDef] (3,2) circle (2.2pt) (3,3) circle (2.2pt) (3,4) circle (2.2pt);
\fill[omDef] (4,3) circle (2.2pt) (4,4) circle (2.2pt);
\draw[omThm,very thick] (0.5,2) -- (4.6,2);
\node[omThm,anchor=west] at (4.8,2) {un instantané~: ce qui était connu en novembre};
\draw[black!60,very thick,dashed] (0.5,4) -- (4.6,4);
\node[black!70,anchor=west] at (4.8,4) {vue d'aujourd'hui~: le dernier millésime};
\draw[omExo,very thick] (2,0.55) -- (2,4.45);
\node[omExo,anchor=west] at (4.8,0.9) {la vie d'un chiffre~: chaque millésime de septembre};
\draw[omExo,->,>=Stealth] (4.75,0.9) -- (2.12,1.4);
\end{tikzpicture}

\omcaption{Données \omterm{def:rs:point-in-time-data-and-the-biases:bitemporal}{bitemporelles}. Chaque point est une valeur stockée pour un mois (\omterm{def:rs:point-in-time-data-and-the-biases:bitemporal}{temps de validité}) telle que publiée à une date (\omterm{def:rs:point-in-time-data-and-the-biases:bitemporal}{temps de connaissance})~; un
mois apparaît pour la première fois le mois suivant et est révisée plus tard. Un backtest doit lire des coupes horizontales à ses instants de décision,
jamais celle du haut~; la ligne verticale passant par septembre représente la vie d'un chiffre, tracée dans la
\cref{fig:rs:point-in-time-data-and-the-biases:life}.}\label{fig:rs:point-in-time-data-and-the-biases:plane}
\end{omfigure}

\section{Biais de survie}

\begin{definition}[Biais de survie]\label{def:rs:point-in-time-data-and-the-biases:survivorship}
Le \emph{biais de survie}\index{biais de survie} (survivorship bias) est l'erreur qui consiste à n'étudier que les titres, fonds ou sociétés qui existent à
la fin de l'échantillon~: ceux qui ont été radiés, liquidés ou absorbés manquent, et leur absence n'a rien d'aléatoire.
\end{definition}

Un univers d'actions construit à partir de la liste actuelle des sociétés cotées a déjà écarté toutes les sociétés qui ont échoué. Une base de
fonds constituée à partir des fonds qui publient aujourd'hui a écarté ceux qui ont fermé. Les valeurs écartées sont, en moyenne, celles qui
ont mal fait, si bien que leur absence embellit toutes les moyennes. Brown, Goetzmann, Ibbotson et Ross (1992) ont montré comment le seul \omterm{def:rs:point-in-time-data-and-the-biases:survivorship}{biais de survie} peut
créer une persistance apparente des performances des fonds~; pour les actions, Shumway (1997) a constaté que les rendements de radiation
des actions radiées pour mauvaise performance étaient pour la plupart absents de la base américaine standard, et que, là où il a pu les reconstituer
à partir des cours de gré à gré, ils s'établissaient en moyenne à $-30~\%$.

\begin{proposition}[Ampleur de la surestimation par les survivants]\label{prop:rs:point-in-time-data-and-the-biases:survivors}
On regroupe toutes les observations titre-mois d'un univers. Soit une fraction $w$ d'entre elles appartenant à des titres radiés avant la
fin de l'échantillon, de rendement moyen $r_F$, et le reste à des titres encore cotés en fin de période, de rendement moyen $r_S$. Le rendement moyen regroupé
de l'univers entier est $(1 - w)r_S + wr_F$, si bien qu'une moyenne calculée sur les survivants le surestime de $w(r_S - r_F)$.
\end{proposition}

\begin{proof}
La moyenne regroupée est la moyenne pondérée des moyennes des deux groupes, les poids étant égaux à leurs parts d'observations titre-mois.
\end{proof}

L'identité indique d'où vient le biais~: non seulement le rendement de radiation, qui ne couvre qu'un seul mois, mais chaque mois pendant lequel un titre défaillant
décline jusqu'à sa radiation. Le marché simulé du chapitre le rend concret. Deux mille actions suivent des marches aléatoires
en logarithme du prix, avec une dérive de 8~\% et une volatilité de 40~\% par an~; une action dont le rendement logarithmique cumulé passe sous $-1.6$ (une perte de
80~\%) est radiée avec un rendement de $-30~\%$ ce mois-là et remplacée par une nouvelle admission à la cote. Sur vingt ans, 2.2~\% des titres sont radiés
chaque année. L'univers à pondération égale, radiations comprises, rapporte 7.7~\% par an~; l'univers des survivants d'aujourd'hui,
complété rétroactivement sur leurs propres historiques, rapporte 12.9~\%~: un biais de 5.2 points par an
(\cref{fig:rs:point-in-time-data-and-the-biases:surv}). Dans les termes de la proposition, 20.9~\% des observations titre-mois appartiennent à des titres
qui échouent plus tard, avec $r_F = -11.1~\%$ par an contre $r_S = 12.6~\%$, soit $w(r_S - r_F) = 5.0$ points~; l'écart restant
vient de la moyenne calculée mois par mois sur un nombre variable de survivants, et non du regroupement. Fixer le rendement de radiation
à zéro abaisse le biais seulement à 4.5 points~: le déclin précédant la radiation en constitue l'essentiel.

\begin{omfigure}
\begin{tikzpicture}
\begin{semilogyaxis}[width=10cm,height=5.2cm,xlabel={années},ylabel={croissance de 1 (échelle logarithmique)},
  tick label style={font=\small},label style={font=\small},xmin=0,xmax=20,ymin=0.8,ymax=16,grid=major,grid style={gray!25},
  ytick={1,2,4,8,16},yticklabels={1,2,4,8,16},legend columns=2,legend style={font=\footnotesize,at={(0.5,-0.32)},
  anchor=north,draw=none,fill=none,/tikz/every even column/.append style={column sep=8pt}},table/col sep=comma]
\addplot[omDef,very thick] table[x=year,y=full]{figdata/research/03-point-in-time-data-and-the-biases/surv.csv};
\addplot[omThm,very thick,dashed] table[x=year,y=survivors]{figdata/research/03-point-in-time-data-and-the-biases/surv.csv};
\legend{tous les titres cotés (avec radiations),survivants d'aujourd'hui seulement}
\end{semilogyaxis}
\end{tikzpicture}

\omcaption{Croissance à pondération égale d'une unité dans un marché simulé de 2\,000 actions sur vingt ans, avec 2.2~\% des titres
radiés chaque année à $-30~\%$~: l'univers investissable finit à 4.6, l'univers des titres encore cotés en fin de période à 12.9.
Données~: le module du chapitre, avec graine fixée.}\label{fig:rs:point-in-time-data-and-the-biases:surv}
\end{omfigure}

\section{Révisions, retraitements et millésimes}

\begin{definition}[Retraitement, millésime de données, biais de remplissage rétroactif]\label{def:rs:point-in-time-data-and-the-biases:vintage}
Un \emph{retraitement}\index{retraitement} est une modification publiée ultérieurement d'une valeur déjà diffusée~: une entreprise qui retraite les comptes
d'un trimestre passé, un institut statistique qui révise un chiffre. Un \emph{millésime de données}\index{millésime de données} est l'ensemble
de l'historique d'une série tel qu'il a été publié à un instant de connaissance donné. Le \emph{biais de remplissage rétroactif}\index{biais de remplissage rétroactif} (backfill bias) apparaît lorsque
de l'historique est ajouté à une base de données après coup, typiquement lorsqu'un fonds ou la nouvelle couverture d'un fournisseur entre avec ses rendements passés~: les entrants
ont choisi d'entrer parce que leur passé était bon, et l'historique rétrorempli n'a jamais été disponible en temps réel.
\end{definition}

La série américaine des emplois non agricoles montre à quel point les révisions pèsent sur un chiffre phare. Le Real-Time Data Set for
Macroeconomists de la Philadelphia Fed (Croushore et Stark, 2001) conserve chaque millésime de la série depuis 1964~; le millésime $m$ est l'historique
tel qu'il figure dans le rapport du mois $m$. Pour les 318 mois de janvier 2000 à juin 2026 disposant d'une première publication, la variation mensuelle
a été révisée ultérieurement de 74\,000 emplois en moyenne en valeur absolue, et de plus de 100\,000 dans 24~\% des mois~; les
révisions sont en moyenne presque nulles ($-3\,000$) mais leur écart type vaut 106\,000. Dix-huit mois ont changé de signe. La plus forte
révision concerne mars 2020 ($-701\,000$ en première publication, $-1\,398\,000$ aujourd'hui).

La donnée de septembre 2008 retrace la vie d'un seul chiffre (la \cref{fig:rs:point-in-time-data-and-the-biases:life})~: $-159\,000$,
puis $-284\,000$ et $-403\,000$ dans les deux rapports suivants, $-321\,000$ après la révision de référence de février 2009, $-458\,000$
après celle de février 2010, $-434\,000$ après celle de février 2011, et $-451\,000$ aujourd'hui. La récession, dans son ensemble, était invisible dans sa
vraie ampleur pendant qu'elle se produisait (la \cref{fig:rs:point-in-time-data-and-the-biases:y2008})~: les premières publications des douze mois de 2008
totalisent une perte de 1.88 million d'emplois, le rapport de janvier 2009 l'évaluait à 2.59 million, et les données actuelles donnent 3.55 million.
Un modèle calibré sur les données d'aujourd'hui apprend une relation entre l'actualité économique et les marchés sur laquelle aucun intervenant n'aurait pu agir.

\begin{omfigure}
\begin{tikzpicture}
\begin{axis}[width=10cm,height=4.8cm,xlabel={mois après septembre 2008 (instant de connaissance)},ylabel={milliers d'emplois},
  tick label style={font=\small},label style={font=\small},xmin=0,xmax=37,ymin=-500,ymax=0,grid=major,grid style={gray!25},
  table/col sep=comma]
\addplot[omDef,very thick,const plot,mark=*,mark size=1pt] table[x=after,y=change]{figdata/research/03-point-in-time-data-and-the-biases/life.csv};
\addplot[black!60,thick,dashed,sharp plot] coordinates {(0,-451) (37,-451)};
\node[black!70,anchor=south,font=\footnotesize] at (axis cs:22.5,-446) {aujourd'hui~: $-451$};
\end{axis}
\end{tikzpicture}

\omcaption{Variation mensuelle de l'emploi américain pour septembre 2008 dans chacun des 36 millésimes mensuels qui suivent~: $-159$, $-284$, $-403$, puis les
révisions de référence de février 2009 ($-321$), février 2010 ($-458$) et février 2011 ($-434$)~; $-451$ dans le millésime d'aujourd'hui
(en pointillés). Données~: Federal Reserve Bank of Philadelphia, Real-Time Data Set for Macroeconomists (données du BLS).}\label{fig:rs:point-in-time-data-and-the-biases:life}
\end{omfigure}

\begin{omfigure}
\begin{tikzpicture}
\begin{axis}[width=11cm,height=5cm,xlabel={mois},xlabel style={yshift=-5pt},ylabel={milliers d'emplois},tick label style={font=\small},
  label style={font=\small},xmin=-0.5,xmax=23.5,ymin=-1000,ymax=200,grid=major,grid style={gray!25},
  ytick={-1000,-500,0},yticklabels={$-1\,000$,$-500$,$0$},
  xtick={0,3,6,9,12,15,18,21},xticklabels={{janv.\\2008},{avr.\\2008},{juil.\\2008},{oct.\\2008},{janv.\\2009},{avr.\\2009},{juil.\\2009},{oct.\\2009}},xticklabel style={align=center},
  legend columns=2,legend style={font=\footnotesize,at={(0.5,-0.42)},anchor=north,draw=none,fill=none,
  /tikz/every even column/.append style={column sep=8pt}},table/col sep=comma]
\addplot[omThm,very thick,mark=*,mark size=1.2pt] table[x=i,y=first]{figdata/research/03-point-in-time-data-and-the-biases/y2008.csv};
\addplot[omDef,very thick,mark=square*,mark size=1.2pt] table[x=i,y=latest]{figdata/research/03-point-in-time-data-and-the-biases/y2008.csv};
\legend{première publication,millésime d'aujourd'hui}
\end{axis}
\end{tikzpicture}

\omcaption{Variations mensuelles de l'emploi américain en 2008--2009 telles que publiées à l'origine et telles qu'on les connaît aujourd'hui. Les premières publications ont sous-estimé les
pertes de 2008 de 1.7 million d'emplois au total. Données~: Federal Reserve Bank of Philadelphia, Real-Time Data Set for
Macroeconomists (données du BLS).}\label{fig:rs:point-in-time-data-and-the-biases:y2008}
\end{omfigure}

Les données d'entreprise se comportent de la même façon. Les comptes sont retraités~; les fournisseurs de données corrigent les erreurs, modifient les définitions et étendent leur
couverture vers le passé. Un fournisseur qui ajoute un millier de petites sociétés dotées de dix ans d'historique crée un univers
rétrorempli, et une stratégie backtestée sur cet univers négocie des sociétés que la firme n'aurait pas pu savoir qu'elle allait couvrir. Le remède est le même dans tous les cas~:
conserver les millésimes.

\section{Jointures as-of et stockage bitemporel}

\begin{definition}[Jointure as-of]\label{def:rs:point-in-time-data-and-the-biases:asof}
Une \emph{jointure as-of}\index{jointure as-of} associe chaque ligne d'une table de gauche (instants de décision) à la dernière ligne d'une table de droite
(observations) dont l'instant est antérieur ou égal à celui de la ligne de gauche, éventuellement à une tolérance près~; elle est stricte
lorsque les lignes de la table de droite ne deviennent connues qu'après leur horodatage.
\end{definition}

Une jointure ordinaire sur des horodatages égaux échoue dans les deux sens~: elle élimine les instants de décision auxquels aucune observation
n'a exactement le même horodatage, et elle associe volontiers une décision à une observation datée du même instant mais arrivée plus tard.
La \omterm{def:rs:point-in-time-data-and-the-biases:asof}{jointure as-of}, fondée sur les instants de connaissance, répond à la seule question qu'un backtest puisse poser~: quelle est la dernière information connue~?

\begin{method}[Rendre un backtest point-in-time]\label{met:rs:point-in-time-data-and-the-biases:pit}
\begin{enumerate}
\item Donner à chaque entrée un instant de connaissance~: l'instant de publication ou de dépôt, l'instant d'arrivée d'un message, la fin de la barre plus le
délai de traitement. Lorsque seul un \omterm{def:rs:point-in-time-data-and-the-biases:bitemporal}{temps de validité} existe, dériver un instant de connaissance prudent d'un décalage documenté.
\item Conserver chaque millésime (données \omterm{def:rs:point-in-time-data-and-the-biases:bitemporal}{bitemporelles})~; ne jamais écraser une valeur.
\item Lire les entrées par l'intermédiaire d'un garde qui connaît l'instant de décision et refuse tout instant de connaissance postérieur.
\item Construire aussi les univers, les identifiants et les ajustements à l'instant de décision (chapitre~4).
\item Tester les fuites d'information~: perturber chaque valeur connue après l'instant de décision et vérifier qu'aucune décision ne change.
\end{enumerate}
\end{method}

\section{Catalogue des autres biais}

Le look-ahead, le \omterm{def:rs:point-in-time-data-and-the-biases:survivorship}{biais de survie} et le remplissage rétroactif sont les biais des données~; les autres proviennent du chercheur. \emph{Sélection de la
période d'échantillon}~: commencer un test après un krach, ou le terminer avant un, est un essai que le journal doit compter (chapitre~1). \emph{Calendriers
désalignés}~: une clôture américaine et une clôture de Tokyo datées du même jour sont séparées de quinze heures. \emph{Cours périmés}~: la dernière transaction d'un instrument peu liquide
peut dater de plusieurs jours, si bien qu'un portefeuille valorisé avec elle paraît plus lisse qu'il ne l'est (chapitre~22). \emph{Opérations sur titres
appliquées avec le recul}~: une série de cours ajustée d'une division avec le facteur d'aujourd'hui transporte dans le passé l'information sur la division
dans le passé (chapitre~4). \emph{Appartenance à un indice}~: tester une stratégie sur les membres actuels d'un indice revient à un \omterm{def:rs:point-in-time-data-and-the-biases:survivorship}{biais de survie}
sous un autre nom. Tous relèvent du look-ahead au sens large, et tous se préviennent par la même discipline~: à chaque instant de décision, seulement
ce qui était connu alors.

\section{Tutoriel~: les chiffres de l'emploi tels qu'ils étaient connus}\label{tut:rs:point-in-time-data-and-the-biases:payrolls}

\begin{tutorial}
\textbf{Objectif.} Charger les millésimes de l'emploi de 2008--2009 dans un entrepôt bitemporel, demander ce qui était connu à quel moment, et laisser un garde refuser
le futur. \textbf{État final~:} les \cref{fig:rs:point-in-time-data-and-the-biases:life,fig:rs:point-in-time-data-and-the-biases:y2008}~;
le total 2008 tel qu'il était connu en janvier 2009 ($-2.59$ million) et aujourd'hui ($-3.55$ million).
\begin{tutsteps}
\item \textbf{L'entrepôt.} Chaque (entité, champ, \omterm{def:rs:point-in-time-data-and-the-biases:bitemporal}{temps de validité}) conserve une liste triée de (\omterm{def:rs:point-in-time-data-and-the-biases:bitemporal}{temps de connaissance}, valeur)~; une requête trouve le
dernier \omterm{def:rs:point-in-time-data-and-the-biases:bitemporal}{temps de connaissance} antérieur ou égal à celui demandé.
\omcode{code/firm/pit/firm_pit.py}{39}{56}{Enregistrer et lire une valeur selon le temps de validité et le temps de connaissance.}\label{lst:rs:point-in-time-data-and-the-biases:store}
\item \textbf{Le garde.} Une vue de l'entrepôt pour un instant de décision~; toute requête portant sur une connaissance postérieure lève une erreur.
\omcode{code/firm/pit/firm_pit.py}{76}{90}{Le garde contre le look-ahead.}\label{lst:rs:point-in-time-data-and-the-biases:guard}
\item \textbf{Interroger.} \texttt{real\_time\_view} pour septembre 2008 avec l'instant de décision d'octobre 2008 renvoie $-159$, et avec
l'instant de décision de septembre 2011, $-434$~; un garde avec l'instant de décision de novembre 2008 lève \texttt{LookAheadError} lorsqu'on demande la connaissance de
décembre 2008. Additionner les douze mois de 2008 tels que connus dans chaque millésime, puis tracer le résultat en regard de la date du millésime.
\end{tutsteps}
\textbf{Pour aller plus loin.} Reconstruire les données depuis la source avec \texttt{rs\_fetch\_payrolls.py}~; ajouter le taux de chômage,
dont le même jeu de données fournit les millésimes, et repérer les mois dont le signe de la variation s'est inversé.
\end{tutorial}

\section{Construction~: le magasin point-in-time}\label{bld:rs:point-in-time-data-and-the-biases:pit}

\begin{build}
\bfield{Rôle} Chaque fait que la petite entreprise utilise en recherche est lu à travers ce magasin, afin qu'un backtest ne voie que ce
qui était connu à chaque instant de décision. Le chapitre~4 construit le référentiel des titres sur cette base, le chapitre~11 les variables fondamentales.
\bfield{Interface} \texttt{Store.put(entity, field, valid, known, value)}~; \texttt{asof(entity, field, valid, known)}~;
\texttt{first}, \texttt{vintages}, \texttt{snapshot(entity, field, known)}, \texttt{latest}~; \texttt{Guard(store,
decision\_time)} avec les mêmes requêtes~; \texttt{asof\_join(left\_t, right\_t, right\_v, tolerance, strict)}.
\bfield{Règles} Ajout seul selon le \omterm{def:rs:point-in-time-data-and-the-biases:bitemporal}{temps de connaissance} (une correction à un \omterm{def:rs:point-in-time-data-and-the-biases:bitemporal}{temps de connaissance} existant le remplace)~; des temps de tout type comparable,
utilisés de manière cohérente~; le garde lève une erreur plutôt que de renvoyer \texttt{None}, afin qu'une fuite soit bruyante.
\bfield{Tests d'acceptation} \texttt{code/firm/pit/tests/}~: la vie du chiffre de septembre 2008~; un instantané égale un millésime et
\texttt{latest} égale la vue d'aujourd'hui~; le garde refuse un \omterm{def:rs:point-in-time-data-and-the-biases:bitemporal}{temps de connaissance} postérieur~; la \omterm{def:rs:point-in-time-data-and-the-biases:asof}{jointure as-of} avec et sans mode strict et
tolérance.
\bfield{Pour aller plus loin} Une version en colonnes indexée par (entité, \omterm{def:rs:point-in-time-data-and-the-biases:bitemporal}{temps de validité}) pour des millions de lignes~; stocker à part le \omterm{def:rs:point-in-time-data-and-the-biases:bitemporal}{temps de connaissance} d'une
correction et celui de la valeur d'origine~; un rapport d'audit de chaque requête faite par un backtest.
\end{build}

\begin{omsources}
\item D.~Croushore et T.~Stark, ``A real-time data set for macroeconomists'', \emph{Journal of Econometrics} 105(1), 2001~;
données~: Federal Reserve Bank of Philadelphia, Real-Time Data Set for Macroeconomists, nonfarm payroll employment.
\item T.~Shumway, ``The delisting bias in CRSP data'', \emph{Journal of Finance} 52(1), 1997.
\item S.~J.~Brown, W.~Goetzmann, R.~G.~Ibbotson et S.~A.~Ross, ``Survivorship bias in performance studies'', \emph{Review of
Financial Studies} 5(4), 1992.
\item E.~J.~Elton, M.~J.~Gruber et C.~R.~Blake, ``Survivor bias and mutual fund performance'', \emph{Review of Financial Studies}
9(4), 1996.
\item W.~Fung et D.~A.~Hsieh, ``Performance characteristics of hedge funds and commodity funds: natural vs.\ spurious biases'',
\emph{Journal of Financial and Quantitative Analysis} 35(3), 2000.
\item R.~W.~Banz et W.~J.~Breen, ``Sample-dependent results using accounting and market data: some evidence'', \emph{Journal of
Finance} 41(4), 1986.
\end{omsources}

\section{Exercices}

\begin{exercise}[$\star$]\label{exo:rs:point-in-time-data-and-the-biases:1}
Une barre quotidienne porte sa date et contient le cours de clôture. Un signal utilise le cours de clôture du jour $t$ et le backtest exécute des transactions à la clôture
du jour $t$. Qu'est-ce qui ne va pas, et quelle est la transaction la plus précoce que le signal autorise~?
\end{exercise}

\begin{exercise}[$\star$]\label{exo:rs:point-in-time-data-and-the-biases:2}
Dans un univers, 15~\% des observations mois-titre appartiennent à des titres radiés par la suite, avec un rendement moyen de $-20~\%$ par an, et les
autres ont un rendement moyen de 10~\%. De combien une moyenne calculée sur les seuls survivants surestime-t-elle celle de l'univers~?
\end{exercise}

\begin{exercise}[$\star$]\label{exo:rs:point-in-time-data-and-the-biases:3}
En utilisant les millésimes de la \cref{fig:rs:point-in-time-data-and-the-biases:life}, que savait un desk sur septembre 2008 en
janvier 2009, en juin 2009 et en septembre 2011~?
\end{exercise}

\begin{exercise}[$\star\star$]\label{exo:rs:point-in-time-data-and-the-biases:4}
Les instants de décision sont 10:00~:00.0, 10:00~:01.0 et 10:00~:02.0~; les cotations sont horodatées 09:59~:59.5, 10:00~:01.0 et 10:00~:01.8, et chacune
arrive 0.3 seconde après son horodatage. Quelle cotation une \omterm{def:rs:point-in-time-data-and-the-biases:asof}{jointure as-of} sur les horodatages attribue-t-elle à chaque décision, et laquelle sur les
instants d'arrivée~?
\end{exercise}

\begin{exercise}[$\star\star$]\label{exo:rs:point-in-time-data-and-the-biases:5}
Les révisions des chiffres de l'emploi, de la première publication à aujourd'hui, ont un écart type de 106\,000 et une moyenne de $-3\,000$. Un modèle de surprise
lit la première publication par rapport à une prévision dont l'erreur, dans les données finales, a un écart type de 90\,000. Quelle part de la
variance apparente des surprises de première publication est du bruit de révision, si les deux sont indépendants~?
\end{exercise}

\begin{exercise}[$\star\star$]\label{exo:rs:point-in-time-data-and-the-biases:6}
Une base de données de hedge funds ajoute des fonds avec leurs trois années de rendements passés lorsqu'ils y entrent. Les fonds entrent lorsque leurs trois années passées
dépassent la médiane. Expliquer le biais dans le rendement moyen de la base et comment le supprimer.
\end{exercise}

\begin{exercise}[$\star\star\star$]\label{exo:rs:point-in-time-data-and-the-biases:7}
\emph{Programmation.} Avec \texttt{survivorship()}, mesurer le biais pour des barrières de $-1.2$ et $-2.0$ et pour des rendements de radiation de $0$
et de $-100~\%$. Quel paramètre déplace le plus le biais, et pourquoi~?
\end{exercise}

\begin{exercise}[$\star\star\star$]\label{exo:rs:point-in-time-data-and-the-biases:8}
\emph{Trouver la faille.} «~Nous avons backtesté notre stratégie de value sur les membres actuels du S\&P~500 depuis 2000, en utilisant les fondamentaux de chaque entreprise
tels qu'ils figurent aujourd'hui dans la base de notre fournisseur. Elle bat l'indice de 3~\% par an.~»
\end{exercise}

\section{Problème~: la prime des survivants}

\begin{problem}[{Problème du week-end --- dans quelle mesure une liste des actions d'aujourd'hui embellit-elle le passé~?}]\label{pb:rs:point-in-time-data-and-the-biases:1}
Le marché simulé du chapitre~: 2\,000 titres, des marches aléatoires sur le log du cours avec une dérive de 8~\% et une volatilité de 40~\% par an, une radiation
lorsque le log du rendement cumulé atteint $-1.6$ avec un rendement de $-30~\%$, un remplacement par une nouvelle cotation, vingt ans, graine 11.

\medskip\noindent\textbf{Partie I --- Le marché.}
\begin{enumerate}
\item Quelle perte de cours représente un log du rendement cumulé de $-1.6$~?
\item Quelle fraction des titres est radiée chaque année~?
\item Quel est le rendement annuel pondéré également de l'univers investissable, radiations comprises~?
\item Et celui des survivants d'aujourd'hui, sur leur propre historique~?
\item Quel est le biais, en points par an~?
\end{enumerate}

\medskip\noindent\textbf{Partie II --- D'où il vient.}
\begin{enumerate}[resume]
\item Quelle part des observations mois-titre appartient à des titres radiés par la suite~?
\item Que valent $r_F$ et $r_S$~?
\item Que donne la \cref{prop:rs:point-in-time-data-and-the-biases:survivors}, et pourquoi diffère-t-elle de la question~5~?
\item Quel est le biais avec un rendement de radiation nul~?
\item Et avec un rendement de radiation de $-100~\%$~?
\end{enumerate}

\medskip\noindent\textbf{Partie III --- Sensibilité.}
\begin{enumerate}[resume]
\item Quels sont le taux de radiation et le biais avec une barrière de $-1.2$~?
\item Et avec une barrière de $-2.0$~?
\item Pourquoi une barrière plus élevée (radiation plus précoce) augmente-t-elle le biais~?
\item Que donnerait un backtest sur les seuls survivants d'une stratégie qui achète des actions après de fortes baisses~?
\end{enumerate}

\medskip\noindent\textbf{Partie IV --- Jugement.}
\begin{enumerate}[resume]
\item Où trouve-t-on les titres radiés pour un backtest réel~?
\item Quel rendement de radiation retient-on lorsqu'une base n'en fournit aucun, et pourquoi pas zéro~?
\item Comment le biais interagit-il avec un univers pondéré également, par rapport à un univers pondéré par la capitalisation~?
\item Énoncer le \emph{résultat nommé}~: la surestimation annuelle du rendement de l'univers des seuls survivants au taux de radiation simulé,
et la part de celle-ci qu'explique le mois de radiation.
\item Nommer deux autres biais de même structure.
\item En une phrase~: pourquoi une liste des entreprises d'aujourd'hui est-elle une prévision~?
\end{enumerate}
\end{problem}

\section{Questions d'entretien}

\begin{interviewq}[$\star$ \iqroles{chercheur, trader}]\label{iq:rs:point-in-time-data-and-the-biases:1}
Qu'est-ce que le \omterm{def:rs:point-in-time-data-and-the-biases:survivorship}{biais de survie}, et comment apparaîtrait-il dans un backtest de sélection d'actions~?
\end{interviewq}

\begin{interviewq}[$\star\star$ \iqroles{chercheur}]\label{iq:rs:point-in-time-data-and-the-biases:2}
Votre modèle macro utilise la croissance du PIB. Quel chiffre de PIB doit-il utiliser pour chaque date du backtest~?
\end{interviewq}

\begin{interviewq}[$\star\star$ \iqroles{développeur, chercheur}]\label{iq:rs:point-in-time-data-and-the-biases:3}
Concevez une table de données fondamentales qui permette des requêtes point-in-time. Quelles sont ses clés~?
\end{interviewq}

\begin{interviewq}[$\star\star$ \iqroles{développeur, ingénieur ML}]\label{iq:rs:point-in-time-data-and-the-biases:4}
Comment vérifieriez-vous automatiquement qu'un pipeline de variables ne comporte aucun \omterm{def:rs:point-in-time-data-and-the-biases:lookahead}{biais d'anticipation}~?
\end{interviewq}

\begin{interviewq}[$\star\star$ \iqroles{chercheur, risque}]\label{iq:rs:point-in-time-data-and-the-biases:5}
Un indice de hedge funds affiche des rendements élevés et une faible volatilité. Quels biais de bases de données soupçonneriez-vous~?
\end{interviewq}

\begin{interviewq}[$\star\star\star$ \iqroles{chercheur}]\label{iq:rs:point-in-time-data-and-the-biases:6}
Montrer qu'une moyenne calculée sur les seuls survivants surestime le rendement moyen d'un univers de $w(r_S - r_F)$, et dire lequel des termes est le plus grand
en pratique.
\end{interviewq}
